\section{Results and discussion} \label{sec:experiments} We report the ceiling and the ceiling-to-model gap,
then the design consequences, then the controls that decide what the gap may be attributed to.
\subsection{Model roster}
\label{sec:roster}
Main text names every model arm by role rather than by vendor identifier (Section~\ref{sec:method}); this
subsection binds each role to its identifier, run conditions, and the accuracies attached to it elsewhere in
the paper, and is the reference for contributions (i) and (ii).
Table~\ref{tab:roster} gives the binding: for each role, the exact identifier, the condition(s) it was run
under, its sample and decode budget, and the accuracies attached to it elsewhere in the paper. The
supervised reference arm discussed alongside \rlFinetuned{} shares the same base model but differs in
training recipe, and is bound on that basis rather than treated as a separate vendor model.

\begin{table*}[t]
\centering
\caption{Role-label roster. Every accuracy figure here is reported elsewhere in the paper under the same
role label; this table adds only the vendor identifier and the run conditions. \rlLadderShort{} and
\rlWeakestShort{} are the short forms used in the columns of Table~\ref{tab:rungs}.}
\label{tab:roster}
\small
\begin{tabular*}{\textwidth}{@{\extracolsep{\fill}}>{\raggedright\arraybackslash}p{0.17\textwidth}>{\raggedright\arraybackslash}p{0.21\textwidth}>{\raggedright\arraybackslash}p{0.27\textwidth}>{\raggedright\arraybackslash}p{0.25\textwidth}@{}}
\toprule
Role label & Identifier & Conditions and sample & Accuracies \\
\midrule
\rlLadder{} (\rlLadderShort) & \ident{gemini-\allowbreak 3.6-\allowbreak flash} & Ladder ($\Rzero$/$\Rone$/$\Rthree$/$\Rfour$), $n{=}210$, $K{=}3$; also the newer arm in the one-generation comparison (same model, same label, not a second arm) &
$\Rzero{=}1.0000$; $\Rone{=}1.0000$; $\Rthree{=}0.8524$; $\Rfour{=}0.7333$; closes 54.1\% of its generational headroom, computed against $\RoneRel$ \\
\rlWeakest{} (\rlWeakestShort) & \ident{gpt-\allowbreak 4.1-\allowbreak mini} & Coordinate condition ($\Rthree$/$\Rfour$), $n{=}210$, $K{=}3$ &
$\Rthree{=}0.4143$ (minimum of the 14 scored); $\Rfour{=}0.3667$ \\
\rlFinetuned{} & \ident{Qwen/\allowbreak Qwen3-\allowbreak VL-\allowbreak 8B-\allowbreak Instruct} & Supervised-then-GRPO, single seed; $n{=}210$; headline at $K{=}8$, also reported at $K{=}3$ &
$0.6905$ at $K{=}8$; $139/210{=}0.6619$ at $K{=}3$; model side of the ceiling-to-model gap (discordance $61{:}6$, $p{=}1.5\times10^{-12}$) \\
Supervised reference arm & \ident{Qwen/\allowbreak Qwen3-\allowbreak VL-\allowbreak 8B-\allowbreak Instruct} (same base as \rlFinetuned{}, supervised-only recipe, no GRPO) & Supervised LoRA, three seeds, 416\,px, 1{,}610 training examples, no augmentation &
Mean $0.6143$ (SD $0.0515$); per-seed $0.590 / 0.567 / 0.686$ \\
\rlBudget{} & \ident{claude-\allowbreak opus-\allowbreak 5} & Two reasoning-budget settings (minimal, provider default), $n{=}210$, $K{=}3$ &
$130/210{=}0.6190$ under both settings; $169/210$ agreement, discordance split $13{:}13$, $p{=}1.0000$ \\
\rlOldGen{} & \ident{gemini-\allowbreak 2.5-\allowbreak flash} & Zero-shot leaderboard, $n{=}210$, $K{=}3$ &
Leaderboard best row $0.4429$; closes 64.0\% of its generational headroom, computed against $\RoneRel$ \\
\bottomrule
\end{tabular*}
\end{table*}

\textbf{Roster composition by condition.} The zero-shot leaderboard scored 13 of a pre-registered 14; the
image-removal control scored 13 of 16; the coordinate condition scored 14 of 15. Every exclusion is a
pre-registered API-error, parse or endpoint gate, and every unscored entry is reported unscored rather than
dropped.

\textbf{Protocol detail: composition-exclusion split.} The random seed fixing the composition-exclusion
split of Section~\ref{sec:benchmark} is 23. The 13 elements reserved for evaluation only are Cd, Ce, Hg, In,
Ir, La, Mn, Os, Re, Ru, Tc, Ti and Tl.

The zero-shot leaderboard behind the headline gap is given in Figure~\ref{fig:leaderboard}; the remaining
secondary-axis figures are in Section~S7 of the ESI\dag{}.

\begin{figure*}[t]
\centering
\includegraphics[width=0.92\textwidth]{figures/leaderboard.pdf}
\caption{No scored model reaches half the oracle's 1.0000 on the original sample. Zero-shot crystal-system
accuracy, $n=210$, $K=3$, 13 models across eight vendors. Solid line: the shape-free baseline on this sample,
0.5286; dashed line: seven-way chance, 0.1429. Error bars: measured run-to-run spread of $K{=}3$ majority
voting, a mean of 5.0 of 210 structures across two sweep runs. Every model on this draw falls below the
shape-free baseline; that pattern is specific to this draw and does not hold at $n = 1933$
(Section~\ref{sec:controls}).}
\label{fig:leaderboard}
\end{figure*}

\subsection{Main results}
\label{sec:results}

\textbf{Rosters.} The three conditions score overlapping but non-identical model sets. The 13 scored zero-shot
models ran on the frozen protocol with a verbatim prompt and no per-model tuning; the best reaches
$93/210 = 0.4429$, far below the oracle's 1.0000. A cell-metric random forest, whose feature specification and
reading sensitivity are detailed in Section~S5 of the ESI\dag{}, reaches 0.8952 on the same label,
above every model, which reproduces an established capability rather than contributing one and is not
comparable to composition-only predictors, since it reads cell geometry rather than chemical composition.

The 0.6905 fine-tuned figure is a single-seed supervised-then-GRPO run at $K{=}8$, distinct from the strongest
zero-shot $R_4$ of 0.7333 on the ladder, and it sits inside the supervised reference arm's across-seed spread,
so \emph{we make no claim that it improves on that arm}; it is used only as the model side of the
ceiling-to-model gap, which survives any seed in that spread.

\subsection{Perception versus the post-perception residual}
\label{sec:residual}

\textbf{Perception dominance is a property of the top of the roster, not of the task.} Supplying ground-truth
geometry as text, the CIF-supplied condition introduced in Related Work and used here as prior methodology,
lifts every model substantially, and none reaches the oracle: the best is 0.8524 against 1.0000
(Figure~\ref{fig:ladder}). The
pre-registered decomposition of Definition~\ref{def:ladder} gives a median perception share of 29.0\% across
the 14 models scored under the coordinate condition (bootstrap 95\% CI $[0.149, 0.358]$), with the
post-perception residual the larger of the two for 13 of them ($p = 0.0018$); the two columns of
Table~\ref{tab:rungs} differ in share by a factor of 5.9. The single exception is the model with the
second-highest $R_4$, tied first on $R_3$, which is why studies confined to frontier models see perception
dominate: their sample sits inside the one region of the roster where it does. Because
$\Rone = \Rzero = 1.0000$ these shares need no reference correction; read against $\RoneRel$ the same
decomposition gives 12 of 14 ($p = 0.0129$), so the conclusion holds at either point of the frontier.

The natural strengthening, that the share rises with model strength, is \emph{not} supported: across the
14 models the share and pixel accuracy are unrelated (Spearman $\rho = -0.1588$, $p = 0.5877$), and we report
the negative rather than the ordering. The raw component moves the other way, $R_3 - R_4$ correlating with
$R_4$ at $\rho = -0.6439$, so the strongest pixel performers are those for whom exact geometry adds least
even though the \emph{share} it closes shows no trend.

\begin{figure*}[t]
\centering
\includegraphics[width=0.92\textwidth]{figures/ladder.pdf}
\caption{Exact geometry as text lifts every model and none reaches the oracle, and the share it closes
tracks no trend in pixel accuracy. Marker shape and colour identify the model in both panels (key below).
(a) Per-model $\Rfour$ (pixels, open marker) and $\Rthree$ (exact geometry as text, filled marker) on the
original evaluation sample ($n=210$, $K=3$), rungs as in Definition~\ref{def:ladder}, with the $\Rone$
oracle at 1.0000 and seven-way chance drawn; $\Rtwo$ is omitted (unscored under its pre-registered gate).
(b) Perception share of the deficit against $\Rfour$: 13 of the 14 models fall below the parity line, so the
deficit that survives exact geometry is the larger one for most of the roster, and the share tracks no trend
in pixel accuracy (Spearman $\rho = -0.1588$, $p = 0.5877$). Shares, the median 0.2901 and its bootstrap
95\% CI $[0.1492, 0.3576]$ (band) are computed against the exact ceiling $\Rone = 1.0000$.}
\label{fig:ladder}
\end{figure*}

\textbf{A control pair makes $R_3$ readable.} The same prompt carrying the chemical formula alone collapses
every scored model toward chance, while with full geometry substituted for images every model reaches between
0.4143 and 0.8524, so the lift is the geometry rather than text-mode prompting.

\subsection{The ceiling is exact; the frontier is in extraction tolerance}
\label{sec:frontier}

\textbf{Identifiability is exact.} Enumerating the phantom set over all ten independent view pairs finds
$\Phi_0 = \emptyset$ on every one of the 2160 certified structures under the released five-camera protocol,
and the oracle returns the correct label for all of them, reconstructing positions to
$3.9\times10^{-15}$\,\AA{} with no atom dropped and no phantom accepted. So $\Rone = 1.0000$ on both samples
at any $\delta \le \tau$, and the hypothesis of Theorem~\ref{thm:ident}(c) is met rather than assumed.
Emptiness is a property of the protocol and not of the construction: the same census over camera subsets
leaves $\Phi_0$ nonempty on 15.4\% of two-view subset-structure pairs and 1.9\% of three-view pairs, so the
fifth camera is doing real work.

\textbf{What degrades is the extraction stage.} A practitioner does not have exact centroids, and the merge
tolerance $\delta$ is the parameter that matters. Reading the same oracle across $\delta$ at fixed $\tau$
traces a frontier: 1.0000 at $\delta \le \tau$, 0.9810 at $5\tau$, and
$\RoneRel = 0.9524$ at the released $15\tau$, with $0.8785$ there on the scaled sample. The loss is entirely
over-merging rather than coincidence, since every structure it costs recovers the correct atom count and has
its closest same-species pair at more than eleven times $\delta$.

This relocates the quantity a reader should treat as the ceiling: identifiability from the renders is total,
and what an actual pipeline achieves is a property of its extraction stage. Every ceiling-to-model gap below
is wider under this reading than at the released tolerance, so no comparison depends on which point of the
frontier is chosen.

\subsection{Controls and intermediate verification}
\label{sec:controls}

\textbf{The residual is an upper bound, and we tested it.} It may include long-list handling rather than
symmetry reasoning: regressing per-structure $R_3$ correctness on atom count, 13 of 14 models show a negative
association ($p=0.0018$) that is not symmetry difficulty in disguise, so $S(m)$
bounds a symmetry-reasoning deficit rather than measuring one.

\textbf{A learned extractor fabricates its output.} Substituting a strong model as the extraction stage,
prompted for species and positions with the symmetry question withheld, then having a weaker model answer from
that text, yields exactly the majority stratum's base rate, so downstream accuracy measures nothing. The
informative measurement is the intermediate, scored against ground truth: median recall 0.0000, with 105 of
206 structures having not one emitted atom within tolerance of a real one despite a median of 48 well-formed
atoms each, where a connected-component blob detector reaches 0.400 on the same task. \emph{A two-stage
pipeline whose intermediate cannot be checked will
attribute fabrication to reasoning}, and only a model-free reference distinguishes the two.

\textbf{The models use the image, and have not memorised the answer.} We ran an image-absent control, the
identical prompt carrying the chemical formula and no renders, which is also the control pair that makes
$R_3$ readable in Section~\ref{sec:residual}. Every scored model collapses toward
chance (mean 0.1487 against 0.1429) and the paired difference is significant for 10 of 13, the three nulls
being the three weakest models. This also bounds pretraining
contamination, which composition exclusion does not: a model that had memorised composition-to-symmetry
mappings would beat chance on the formula alone, and none does.

\textbf{Image-removal control, per model.} Figure~\ref{fig:noimage} gives the per-model breakdown: each model answered the identical prompt with the five
renders and with the chemical formula alone, and the paired difference is significant for 10 of the 13
scored, the three nulls being the three weakest models.


\begin{figure*}[t]
\centering
\includegraphics[width=0.92\textwidth]{figures/noimage.pdf}
\caption{Removing the image collapses every model toward chance. Each model answered the identical prompt
with the five renders and with the chemical formula alone, paired per structure (original sample, $n=210$,
$K=3$, 13 scored of 16 attempted); model identifiers are drawn at the end of each bar pair. Stars mark
models whose paired difference is significant, 10 of the 13, the three nulls being the three weakest models
(two of which are at chance even with images); the dashed line is seven-way chance.}
\label{fig:noimage}
\end{figure*}

\textbf{A below-baseline comparison is scoped to its sample, and does not survive scale-up.} On the
original sample every model falls below a shape-free baseline reading only atom count, density and cell
volume (0.4429 against 0.5286, one-sided $p = 0.0078$), but a larger set at $n = 1933$ generated under
identical rules puts that baseline at 0.2054, so \emph{the 0.5286 is a property of the original draw, not of
the task}. The same size-matched control shows the oracle's apparent drop at scale \emph{is} a cell-size
effect, 0.8774 on the full sample against 0.9459 at matched cell size, since more atoms per cell means more
candidate coincidences of the kind Definition~\ref{def:phantom} counts. The ceiling-to-model gap therefore
survives scale-up while the baseline comparison does not.

\textbf{What the controls do and do not isolate.} Three competing explanations for the gap are excluded by
design: that the renders lack the answer, by $R_1$ and the occluder classification; that models pattern-match
or have memorised the label, by the image-removal control; that the $R_3$ lift is text mode rather than
geometry, by the formula-only arm. Two are not excluded: the residual is an upper bound, and the extraction
share is bounded from neither side. A second execution of the
sweep differs by 5.0 structures on average, so adjacent leaderboard rows are not separable.

\subsection{Render conventions}
\label{sec:conventions}

At exact extraction the frozen protocol, the off-axis camera perturbation and the released tiled geometry
each return all 210 labels, so no camera choice in this family withholds any part of the task; what separates
them is conditioning. Before measuring we fixed $\kappa = 2.7321$ for the frozen cameras and
$\kappa = 8.4258$ for the off-axis perturbation, predicted that the more poorly conditioned protocol would
tolerate proportionally less centroid error, and recorded a noise budget from the band of
Proposition~\ref{prop:scope}. Injecting per-view Gaussian centroid noise returns one confirmation and two
failures, and we report all three. The prediction that held is the one that isolates $\kappa$: frozen and
off-axis differ in camera placement alone, and the frozen protocol tolerates $2.1\times$ the noise, with
paired comparison separating them at $\sigma = 0.02$\,\AA{} (102 structures against 2,
$p = 5.4\times10^{-28}$). Two predictions failed, the three-way ordering including the tiled arm and the
absolute noise budget, the latter optimistic by a factor of 31 because the band is sufficient rather than
tight. So $\kappa$ orders protocols differing only in cameras and delivers neither a usable noise budget nor a
cross-density ordering; stating both is the point of fixing the prediction first. View count behaves as
Corollary~S1(a) of the ESI\dag{} says, and what the protocol withholds through occlusion remains under 1\% of
atoms and changes no classification. The noise grid, both failure mechanisms, per-pair $\kappa$ values, the
density decomposition and the probe power analysis are in Section~S6 of the ESI\dag{}.

\textbf{Occluder classification.} This is the occlusion audit referenced in Section~\ref{sec:controls}.
Because the oracle assumes ideal extraction, occluders in the released
renders were classified by whether each one hides a symmetry-equivalent copy of the atom it occludes; most
occlusion hides nothing recoverable, and since triangulation needs only two clear views this holds for
under 1\% of atoms (0.26\% and 0.87\% of the two samples) even before classification. The redundant-only
control removes only occluders in that recoverable class, leaving informative occlusion untouched.
Figure~\ref{fig:conditions} reports recovery under all three conditions on both evaluation samples.

\begin{figure*}[t]
\centering
\includegraphics[width=0.92\textwidth]{figures/conditions.pdf}
\caption{Removing informative occlusion from the oracle's input changes no classification, so what the
renders hide is small and redundant, and the removals do reach the reconstructor. Expansion sample,
$n = 210$, four views. (a) Symmetry recovery under the four occlusion conditions: O0 unconditioned, O1
informative occluders removed, O2 all occluders removed, O3 the pre-registered random-removal control; the
dashed line is O0. (b) Exact atom-count match under the same conditions, which moves across them, so the
flat recovery in (a) is a property of what occlusion hides rather than a failure of the intervention to
reach the reconstructor.}
\label{fig:conditions}
\end{figure*}

\subsection{Limitations and scope}
\label{sec:limits}
 The construction needs a known, invertible
forward rendering, which holds for procedurally generated benchmarks generally but is instantiated here on one
task in one domain, so what is not shown is that a second instantiation reproduces the same qualitative
findings; it does not depend on orthography as heavily as the algebra suggests, since under calibrated
perspective Theorem~\ref{thm:ident} survives while Corollary~S1(b) of the ESI\dag{} fails and $\kappa$ becomes
position-dependent. The oracle assumes ideal extraction, so $R_1$ bounds what is recoverable rather than what
is achievable, and the one extractor evaluated end to end fails the hypothesis of
Proposition~\ref{prop:scope} on soundness, completeness and localisation together; a visibility-conditioned
variant is a distinct measurement we have not built, and no human baseline exists. Because $\Phi_0$ is empty
here, identifiability contributes nothing to the ceiling-to-model gap, which sharpens the instrument but means
the single headline number a reader might want does not exist: what exists is the frontier plus the
class-membership statement of Proposition~\ref{prop:scope}, and $\kappa$ governs one axis of that frontier
only, so it is a design comparator rather than a tolerance budget. Finally, the labels are computed rather
than annotated, but the databases they are cross-checked against derive their own assignments through the same
library (Section~\ref{sec:benchmark}), so certification establishes stability within one algorithm family
rather than agreement between distinct ones; the rosters carry no current-frontier arm beyond the three
compared and no prompt-sensitivity sweep was run, the remaining secondary-axis findings are collected in
Section~S7 of the ESI\dag{}, $R_3$ used a single textual schema, and the model arms were not
re-run at $n=1933$, so that gap is bounded on the oracle side only. Each limit is given in full in
Section~S8 of the ESI\dag{}.

An oracle that recovers structure from a render measures information, not scientific competence, and a
ceiling reported here should not be read as evidence that a model, or a person, can solve the underlying
task. Every ceiling in this article is reported with the tolerance at which it was read, so no number can be
quoted without the regime that produced it, and the per-structure prediction vectors are released so any
ceiling can be re-derived rather than cited.

